\documentclass[10pt,a4paper]{amsart}
\usepackage[margin=19mm]{geometry}
\usepackage{amsmath,amssymb,mathtools}
\usepackage[scr=rsfs,cal=euler]{mathalfa}
\usepackage{xfrac}
\usepackage[unicode,hidelinks,bookmarks=false]{hyperref}
\newcommand{\ie}{i.e.\ }
\newcommand{\cf}{cf.\ }
\newcommand{\resp}{respectively\ }
\newcommand{\Subsection}{\subsection}
\providecommand{\nobreakdash}{\nobreak\hbox{-}}
\setlength{\emergencystretch}{2em}
\multlinegap0pt
\usepackage{amssymb}
\usepackage{tikz}
\newtheorem{theorem}{Theorem}
\newcommand{\calR}{\mathcal{R}}
\newcommand{\op}{\mathrm{op}}
\newcommand{\tmap}{\mathsf{T}}
\title{Equivariant quantizations of the positive nilradical and covariant differential calculi}
\author{Marco Matassa}
\date{Source: arXiv:2404.18544v2 (CC BY 4.0)}
\begin{document}
\maketitle

\noindent\emph{Source and licence.} Equivariant quantizations of the positive nilradical and covariant differential calculi. Version arXiv:2404.18544v2 (CC BY 4.0), distributed by arXiv under the Creative Commons Attribution 4.0 International licence, http://creativecommons.org/licenses/by/4.0/, as recorded in the arXiv licence metadata for this exact version and shown on the abstract page https://arxiv.org/abs/2404.18544v2. Full licence notice: \texttt{licenses/arxiv-2404.18544-CC-BY-4.0.txt}.\par\smallskip

\noindent\textit{Editorial scope.} Counted unit: the article's theorem thm:transmutation (source0.tex lines 1173--1181). The article's complete original proof of it is delivered with it (source0.tex lines 1183--1232, one \texttt{proof} environment of 50 lines). No mathematics is written, repaired or re-derived: the statement and the proof are the article's own lines, in the article's own notation, and no retained line is substituted or re-flowed.  No further source-local context is needed: every cross-reference of every retained line resolves inside the two delivered spans.  Nothing was omitted from any retained span, so the package carries no omission record.  No retained line contains a citation, so the wrapper carries no bibliography and no bibliography entry.  No retained line contains a figure environment, an image-inclusion command or drawing code, so no image is delivered and the package declares no asset dependency.  Editorial formatting only, in addition: an AMS \texttt{amsart} preamble that restates the article's own declarations and macros used by the retained lines, namely the environments \texttt{theorem} on separate counters, together with the standard packages those lines need.  The retained environments print in this document as Theorem 1 in order of appearance, whereas the article prints the same items with the numbering of its own full document.  The complete native source of the article as distributed in the arXiv e-print for this exact version is bundled unchanged as \texttt{sources/158307/source0.tex}.
\begin{theorem}
\label{thm:transmutation}
Let $H$ be a quasitriangular Hopf algebra with universal R-matrix $\calR = \sum_i a_i \otimes b_i$ and braiding $\sigma$.
Define the map $\tmap: H \otimes H \to H \otimes H$ by
\[
\tmap(x \otimes y) = \sum_i x b_i \otimes a_i \triangleright y.
\]
Then we have an isomorphism of $H$-module algebras $\tmap: (H \otimes_\sigma H, \triangleright) \to (H \otimes H, \blacktriangleright)$.
\end{theorem}

\begin{proof}
\textbf{Invertibility}. Using the properties of the R-matrix one easily checks that
\[
\tmap^{-1}(x \otimes y) = \sum_i x S(b_i) \otimes a_i \triangleright y.
\]
\textbf{Equivariance}. Using the definition of the left action $\blacktriangleright$ we compute
\[
\begin{split}
z \blacktriangleright \tmap(x \otimes y)
& = \sum_i z_{(1)} x b_i S(z_{(3)}) \otimes z_{(2)} a_i \triangleright y \\
& = \sum_i z_{(1)} x b_i S(z_{(3)}) \otimes z_{(2)} a_i S(z_{(4)}) z_{(5)} \triangleright y.
\end{split}
\]
Using $\calR_{2 1} \cdot \Delta^\op(x) = \Delta(x) \cdot \calR_{2 1}$ and $\Delta^\op(S(x)) = S(x_{(1)}) \otimes S(x_{(2)})$ we get
\[
\begin{split}
z \blacktriangleright \tmap(x \otimes y)
& = \sum_i z_{(1)} x S(z_{(4)}) b_i \otimes z_{(2)} S(z_{(3)}) a_i z_{(5)} \triangleright y \\
& = \sum_i z_{(1)} x S(z_{(2)}) b_i \otimes a_i z_{(3)} \triangleright y \\
& = \sum_i (z_{(1)} \triangleright x) b_i \otimes a_i \triangleright z_{(2)} \triangleright y = \tmap(z \triangleright (x \otimes y)).
\end{split}
\]
\textbf{Homomorphism}. The coproduct identities for $\calR = \sum_i a_i \otimes b_i$ can be written as
\begin{equation}
\label{eq:identities-R}
\begin{split}
\sum_i a_{i (1)} \otimes a_{i (2)} \otimes b_i & = \sum_{i, j} a_i \otimes a_j \otimes b_i b_j, \\
\sum_i a_i \otimes b_{i (1)} \otimes b_{i (2)} & = \sum_{i, j} a_i a_j \otimes b_j \otimes b_i.
\end{split}
\end{equation}
Using $x y = (x_{(1)} \triangleright y) x_{(2)}$ and the second identity in \eqref{eq:identities-R}, we compute
\[
\begin{split}
\tmap(x \otimes y) \cdot \tmap(x' \otimes y')
& = \sum_{i, k} x b_i x' b_k \otimes (a_i \triangleright y) (a_k \triangleright y') \\
& = \sum_{i, k} x (b_{i (1)} \triangleright x') b_{i (2)} b_k \otimes (a_i \triangleright y) (a_k \triangleright y') \\
& = \sum_{i, j, k} x (b_j \triangleright x') b_i b_k \otimes (a_i \triangleright a_j \triangleright y) (a_k \triangleright y').
\end{split}
\]
On the other hand, using the first identity in \eqref{eq:identities-R} we compute
\[
\begin{split}
\tmap(x \tilde{x} \otimes \tilde{y} y')
& = \sum_i x \tilde{x} b_i \otimes (a_{i (1)} \triangleright \tilde{y}) (a_{i (2)} \triangleright y') \\
& = \sum_{i, k} x \tilde{x} b_i b_k \otimes (a_i \triangleright \tilde{y}) (a_k \triangleright y').
\end{split}
\]
Setting $\tilde{x} \otimes \tilde{y} = \sum_j b_j \triangleright x' \otimes a_j \triangleright y$ gives $\tmap(x \tilde{x} \otimes \tilde{y} y') = \tmap(x \otimes y) \cdot \tmap(x' \otimes y')$.
Finally we note that since $\tilde{x} \otimes \tilde{y} = \sigma(x' \otimes y)$ we have $x \tilde{x} \otimes \tilde{y} y' = (x \otimes y) \cdot_\sigma (x' \otimes y')$.
\end{proof}

\end{document}
